\chapter{序論}

% \section{研究背景}
% ぼつ

% スマートフォンの壁紙は、端末を使用する際に日常的かつ継続的に目にする視覚コンテンツの一つである。

% 一般的な静止画によるスマートフォン壁紙に対し、ライブ壁紙では時間経過に伴う描画の変化に加え、端末のセンサ情報やユーザの操作などを利用することで、インタラクティブな表現を実現することができる。
% Androidでは、アプリケーションから独自のライブ壁紙を実装するためのWallpaperServiceが提供されており、背景上で継続的に動作する動的な表現を構築することが可能である\cite{android_wallpaper_service}。
% このことから、ライブ壁紙は単なる装飾としての背景画像にとどまらず、日常的に利用される端末上で動的・インタラクティブな視覚表現を提示する媒体として捉えることができる。

% ライブ壁紙は、映像作品や一般的なアプリケーションのように、ユーザが能動的にコンテンツを起動して鑑賞するものとは異なり、スマートフォンの日常的な利用の中で継続的に接触することのできる媒体である。

% 端末や設定によってはホーム画面だけでなくロック画面にも表示されるため、アプリケーションを起動することなく視覚表現を提示することができる。

% このような常時性およびアクセスの容易さは、ライブ壁紙を日常生活に近接したインタラクティブ表現媒体として利用できる要因の一つである\cite{dearman2012wallpaper,murray2013avatar,nwagu2023chai}。




% また近年、Androidにおいてライブ壁紙のパーソナライズを支援する機能の拡張も行われている。

% Android 16では、ユーザが提供するコンテンツなどを含む動的なライブ壁紙を扱うため、WallpaperDescriptionおよびWallpaperInstanceが導入された\cite{android16_wallpaper}。
% 同一のライブ壁紙サービスから異なる内容を持つ壁紙を扱うことが可能となるなど、動的かつユーザ主導のライブ壁紙を実現するための基盤が拡張されている。
% このことからも、ライブ壁紙において、単一の固定された表現を提示するだけでなく、ユーザや利用状況に応じて内容を変化させることには検討の余地があると考えられる。

\section{研究背景}

スマートフォンでは、ホーム画面やロック画面の背景として、ユーザが任意の画像を壁紙に設定することができる。壁紙は端末を利用するたびに繰り返し目にするものであり、スマートフォンにおける身近な視覚要素の一つである。

一般的な壁紙には静止画像が用いられるが、スマートフォンには、画像や映像が時間とともに変化する「ライブ壁紙」も存在する。ライブ壁紙には、単にアニメーションを再生するものだけでなく、時刻に応じて景色が変化するもの、画面へのタッチに反応するもの、端末を傾けることで背景の見え方が変化するものなど、さまざまな形式がある。

このようなライブ壁紙では、スマートフォンが備えるタッチ入力や各種センサを利用することで、ユーザの操作や端末の状態に応じて表示を変化させることができる。そのため、ライブ壁紙は単なる装飾としての背景ではなく、スマートフォンの画面全体を利用した動的・インタラクティブな視覚表現の媒体として捉えることができる。

また、ライブ壁紙は、一般的な映像作品のように鑑賞のために再生操作を行う必要がなく、スマートフォンを日常的に利用する過程で自然に目にするという特徴を持つ。この性質を利用し、壁紙を情報提示やユーザ行動へのフィードバックに用いた研究も行われている\cite{dearman2012wallpaper,murray2013avatar,nwagu2023chai}。

近年では、ライブ壁紙に対して、単一の固定された表現を提示するだけでなく、ユーザや利用状況に応じて内容を変化させるための仕組みも拡張されている。Android 16では、ユーザが提供するコンテンツなどを含む動的なライブ壁紙を扱うために、WallpaperDescriptionおよびWallpaperInstanceが導入されており、同一のライブ壁紙サービスから異なる内容を持つ壁紙を扱うことが可能となっている\cite{android16_wallpaper}。

このことから、ライブ壁紙は、あらかじめ定められた一つの表現を提示するだけの媒体ではなく、ユーザごとの好みや利用状況に応じて表示内容を変化させることのできる視覚表現媒体として発展する余地があると考えられる。特に、表現の種類や強さを利用者自身が調整できる仕組みを取り入れることで、より個人の選好に応じたライブ壁紙を設計できる可能性がある。

本研究では、このようなライブ壁紙の表現媒体としての特徴に着目し、その中でもグリッチ表現を対象とする。
グリッチは、本来、デジタル機器やデータ処理などで生じる予期しないエラーや破損を指すが、それらに伴う視覚的・聴覚的な現象を意図的に作品へ取り入れるグリッチアートが展開されてきた。
Menkmanは、グリッチを単なる視覚効果としてではなく、技術的なエラー、その知覚、およびメディアとの関係から論じている\cite{menkman2011momentum}。
また、偶発的に発生したエラーだけでなく、エラーに特徴的な表現を意図的に生成する手法も、グリッチアートの一形態として論じられている。

グリッチ表現は静止画や映像作品のみならず、リアルタイムな表現にも利用されている。
例えば井藤らは、Webカメラから取得した映像に対してリアルタイムにデータモッシングを適用する手法を提案しており、あらかじめ生成された映像を再生するのではなく、入力される映像に応じてその場で変化するグリッチ表現を実現している\cite{ito2013datamoshing}。
このような入力とグリッチ表現との結び付きは、グリッチを固定された視覚効果ではなく、インタラクティブな表現として扱う可能性を示している。

ライブ壁紙はスマートフォンの画面背景として広い領域に表示され、
その上にアイコンやユーザインタフェースが重なるという特徴を持つ。
そのため、独立した映像領域の内部で表現を提示する場合と比較して、
表示内容を端末そのものの状態と結び付けて知覚させる表現が可能であると考えられる。

例えば、映像作品中にグリッチを提示した場合、
ユーザはそれを映像コンテンツ内部で生じている表現として認識する。
これに対して、ホーム画面上でグリッチが継続的に表示され、
さらに端末の傾きや操作に応じてその状態が変化する場合、
グリッチが映像の内部に限定された現象ではなく、使用している端末そのものに生じているような表現として提示できる可能性がある。
さらに、端末の傾きなどの入力に応じて表示を変化させることで、ユーザ自身の操作とグリッチ表現とを直接結び付けることができる。
これは、故障やノイズといったデジタル機器に由来する視覚的特徴を持つグリッチと、端末画面全体を表示領域とするライブ壁紙との組み合わせに固有の表現上の特徴であると考えられる。






\section{研究課題}

グリッチアートでは、デジタル画像や映像にエラーやノイズを生じさせる多様な手法が用いられており、リアルタイムな入力を利用する表現についても研究が行われている。
一方で、これらの表現を日常的に使用されるスマートフォン上で、ユーザの操作に応じて変化するものとして手軽に利用できる環境は限定的である。
特に、本研究で対象とするような、ライブ壁紙とグリッチ表現を組み合わせ、端末操作とのインタラクションを含めて表現設計を検討した例は多くない。

そこで本研究では、当初、グリッチ表現をスマートフォン向けライブ壁紙として実装することにより、日常的かつインタラクティブに利用できるグリッチ表現環境の構築を課題とした。

その後、制作したライブ壁紙を用いた予備評価を行ったところ、グリッチ表現に対する評価や選好には参加者間で大きな差が見られた。
ある表現を好意的に評価する参加者が存在する一方、同じ表現を好まない参加者も存在し、単一の表現をすべてのユーザに対して提示する設計では、ライブ壁紙としての好ましさを十分に満たせない可能性が示唆された。

壁紙はユーザが日常的に目にするものであり、静止画壁紙ではユーザ自身が好みの画像を選択することが一般的である。
一方、動的なライブ壁紙では、動き方や色、効果の強度など、静止画よりも多くの表現要素を含むため、それらに対する選好にも個人差が生じると考えられる。
そこで本研究では、インタラクティブなグリッチ表現を実現することに加え、ユーザが自身の選好に応じて表現を調整できるライブ壁紙について検討することを新たな課題とする。


\section{研究目的}

本研究の目的は、スマートフォン上で日常的に利用可能なインタラクティブグリッチ表現を実現するとともに、ライブ壁紙におけるグリッチ表現に対するユーザの選好を明らかにし、その個人差に対応するための表現設計について検討することである。

具体的には、Android端末を対象として、端末の傾きなどの入力に応じて視覚表現が変化するグリッチライブ壁紙を設計・実装する。
さらに、複数のグリッチ表現を用いた評価を行い、各表現に対するユーザの印象および選好を調査する。
その結果を踏まえ、色やグリッチ効果などをユーザ自身が調整可能なパーソナライズ機能を導入し、ライブ壁紙における個人差を考慮した表現設計のあり方を検討する。

本研究を通じて、特定のグリッチライブ壁紙を制作することだけでなく、動的なライブ壁紙を設計する際にユーザごとの選好を考慮する必要性を示し、今後のインタラクティブなライブ壁紙制作における設計上の知見を得ることを目指す。


\section{本研究のアプローチ}

本研究では、グリッチ表現を生成しAndroidライブ壁紙として表示するシステムを構築する。

まず、生成AIを用いてグリッチ表現の基礎となる画像を生成する。
次に、生成された画像を解析し、画像に含まれる色情報などの特徴を抽出してJSON形式で保持する。
ライブ壁紙の描画時には、この情報を読み込み、元画像に対応した色を用いてグリッチ表現を生成する。
これにより、あらかじめ固定された色のグリッチを重ねるのではなく、元となる画像の特徴を反映した表現を生成する。

Android上の描画にはOpenGL ESを用い、画像の変形、ノイズ、色ずれなどの複数のグリッチ表現をリアルタイムに生成する。
また、端末のセンサから傾き情報を取得し、視差表現やグリッチの変化に反映することで、ユーザの端末操作に応じたインタラクティブなライブ壁紙を実現する。

さらに、制作した複数の表現についてユーザ評価を行い、グリッチの種類や強度等に対する選好を調査する。
予備評価によって確認された選好の個人差を踏まえ、システムに表現を調整するための機能を追加し、パーソナライズ可能なライブ壁紙として改良する。
その後、改良したシステムおよび評価方法を用いて再評価を行い、ライブ壁紙におけるグリッチ表現とパーソナライズについて考察する。


\section{本論文の構成}

本論文は全9章で構成する。

### 本論文の構成

第1章では、本研究の背景、研究課題、目的および研究のアプローチについて述べる。

グリッチアートおよびライブ壁紙の現状を概観し、両者を組み合わせた表現に着目した経緯を示す。
その上で、ライブ壁紙という日常的な利用環境においてグリッチ表現を扱う意義と、本研究が解決を目指す課題を明らかにする。


第2章では、本研究の基礎となる概念および関連研究・関連事例を整理し、本研究の位置付けを示す。

まず、グリッチアートとライブ壁紙の定義および特徴を整理し、関連事例を通して、それぞれの表現や利用形態について検討する。
また、グリッチ表現の種類と特徴を整理し、本研究で用いる各表現の名称および分類を定める。
これらを踏まえ、既存の表現や研究との関係から、本研究が対象とする領域を明確にする。


第3章では、本研究で用いる動的グリッチ表現の設計について述べる。

各グリッチ表現の視覚的特徴、表現を構成する要素、および端末操作との対応関係を示す。
また、制作した第1版と第2版を比較し、表現の構成や視覚的な変化を整理するとともに、それぞれの設計意図と改良の経緯を説明する。


第4章では、提案するグリッチライブ壁紙システムの構成および実装について述べる。

生成AIを利用した元画像の生成、画像からの色情報の抽出、JSON形式による情報の保持、グリッチ生成処理、OpenGL ESによる描画、センサ入力、およびユーザインタフェースについて説明する。
さらに、各処理の連携とデータの流れを示し、提案した表現がライブ壁紙上で実現される仕組みを明らかにする。


第5章では、制作したグリッチライブ壁紙を用いた予備評価について述べる。

評価の目的、対象者、提示した刺激、実施手順および質問項目を示し、得られた結果を整理する。
また、学会発表において得られた意見や指摘を取り上げ、制作物および評価方法に関する課題を明らかにする。


第6章では、予備評価および学会発表から得られた課題を踏まえ、システムと評価方法に加えた改良について述べる。

特に、グリッチ表現の視覚的な改善、ユーザインタフェースの見直し、および利用者による表現の調整機能について説明する。
また、予備評価の限界を整理し、表現に対する選好とその個人差をより適切に捉えるための評価方法を検討する。


第7章では、改良したシステムを用いた評価について述べる。

実験条件、評価対象、実施手順および評価指標を示し、得られた結果を分析する。
特に、グリッチ表現の種類や強さ、端末操作との対応、および利用者による調整が、表現に対する評価にどのように関係するかを検討する。


第8章では、予備評価および改良後の評価から得られた結果を総合し、ライブ壁紙における動的グリッチ表現について考察する。

グリッチらしさ、視覚的な好ましさ、壁紙としての利用意向などの関係を整理し、表現に対する選好の個人差について検討する。
また、端末操作に応じたインタラクションの役割と、利用者による表現の調整が持つ意義を考察し、パーソナライズを取り入れたライブ壁紙の設計に関する知見をまとめる。


第9章では、本研究で得られた知見を総括し、今後の課題について述べる。

特に、端末の性能や画面特性の違いによる影響、長期間の利用に伴う評価の変化、および実際の利用環境における視認性や操作性について検討する。
さらに、これらを踏まえ、動的グリッチ表現をライブ壁紙として継続的に利用するための設計上の課題と、今後の展望を示す。

